\section{预期的成果}

\subsection{成果一：威胁感知的工具调用链风险校准与分级防御策略（TARC-Guard）}

\begin{enumerate}
  \item \textbf{统一日志模式与任务集。}完成面向工具增强型LLM Agent的统一日志模式设计（含\texttt{src}/\texttt{sink}/\texttt{confirm}字段），形成覆盖文件管理、邮件办公、数据库查询、日程管理、代码执行和IoT控制六类工具场景的低成本任务集，包含600条以上工具调用链样例，覆盖$R_0$至$R_6$全部风险类型。

  \item \textbf{七维风险因子体系与校准评分框架。}实现工具敏感度、权限扩散度、调用深度、不可逆操作比例、敏感数据暴露度、人工确认覆盖率和任务必要性偏差（$N_{\text{need}}$）七个可解释因子的自动计算，以及校准风险评分函数$R_\theta(C) = \sigma(b + \theta^\top \Phi(C))$和因子贡献分解输出。通过“威胁先验+弱监督+约束优化”四步整定策略学习权重，并报告ECE和Brier Score验证概率可解释性。

  \item \textbf{L0--L4分级防御策略与防御动作优化。}实现L0--L4五级风险等级划分和防御动作优化公式（式~\ref{eq:guard_action}），输出分级防御建议（Allow/Log/Mask/AskConfirm/Block+Rollback）。通过对比实验（与单步权限检查、敏感工具计数等基线对比）和消融实验（各因子独立贡献分析）验证框架有效性。
\end{enumerate}

\subsection{成果二：基于行为轨迹图的工具滥用威胁归因与自适应防御方法（TraceGraph-Guard）}

\begin{enumerate}
  \item \textbf{带防御节点的异构行为轨迹图建模方法。}提出适用于LLM Agent执行过程的异构行为轨迹图表示，节点集合新增防御节点$V_{\text{guard}}$，边集合新增防御边$E_{\text{defend}}$，将自然语言任务、计划、工具调用、工具参数、返回结果、记忆写入、最终输出和防御动作统一到一个可查询的图结构中，并实现确定性图构建算法。

  \item \textbf{三级升级判定检测机制（TTED）与图规则库。}构建面向A1--A6六类攻击的8条核心图规则（Rule-1至Rule-8），结合类型化路径特征+轻量分类器（GAT可选增强）和LLM终审实现三级检测。图级综合风险评分融合序列先验：$M(G_t) = \lambda_1 S_{\text{rule}} + \lambda_2 S_{\text{cls}} + \lambda_3 S_{\text{llm}} + \lambda_4 R_\theta(C)$。在AgentDojo和AgentCanary基准上完成与代理基线方法的对比实验，验证图检测相对于序列方法的优势。

  \item \textbf{风险路径定位、反事实归因与最小干预防御。}实现含$-\eta \cdot confirm(p)$保护项的Top-K风险路径定位算法、基于反事实消融的节点/边归因方法，以及最小干预防御公式（式~\ref{eq:min_intervention}），输出包含异常类型、风险路径、关键节点、防御动作和恢复计划的完整证据链报告$Defense(G_t)$。通过PathRecall@K、Hit@K、MRR、DSR（防御成功率）和UPR（效用保持率）等指标定量验证归因准确性和防御有效性。
\end{enumerate}

\subsection{成果三：可插拔Defense Guard中间层原型系统}

基于整体研究内容，设计并实现Defense Guard可插拔中间层原型系统，系统包含七大核心模块：

\begin{enumerate}
  \item \textbf{Trace Collector（轨迹采集模块）：}采集或读取Agent执行日志，规范化为统一模式（含\texttt{src}/\texttt{sink}/\texttt{confirm}字段），支持多种Agent框架的日志格式。
  \item \textbf{TARC-Guard Engine（风险校准引擎）：}计算七维风险因子，输出校准风险分数$R_\theta(C)$、因子贡献分解和L0--L4分级防御动作建议。
  \item \textbf{Graph Builder（图构建模块）：}将日志转换为带$V_{\text{guard}}$/$E_{\text{defend}}$的异构行为轨迹图，支持可视化输出。
  \item \textbf{TTED Detector（三级检测模块）：}依次执行图规则拦截、类型化路径特征分类器和LLM终审，输出图级综合风险评分$M(G_t)$。
  \item \textbf{Attribution Analyzer（归因分析模块）：}输出Top-K风险路径、节点/边归因贡献分解，生成可解释归因报告。
  \item \textbf{Defense Guard（防御执行模块）：}根据最小干预防御公式选择防御动作，通过注入\texttt{defense\_constraint}属性实现Pre-call/In-call/Post-call/Evidence Guard四阶段防御。
  \item \textbf{Audit Report Generator（审计报告生成模块）：}生成包含归因路径、防御动作和恢复计划的完整证据链报告$Defense(G_t)$，支持命令行和批量处理。
\end{enumerate}

\subsection{学术论文}

基于研究内容与实验结果，计划撰写并发表会议论文一篇。
